
\section{Programming Paradigms}

A programming paradigm may be understood as a conceptual model for organizing
and expressing computation. Different paradigms provide distinct abstractions
for representing data, state, operations, and control flow
\cite{scott2015programming}.

The two principal paradigms considered in this work are imperative programming
and functional programming.

In the imperative paradigm, a program can be understood as a sequence of
instructions that modify computational state. Assignments, loops, and
conditional structures explicitly determine how the state of the program
changes during execution
\cite{scott2015programming,kernighan1988c}.

For example, a numerical recurrence

\begin{equation}
    x_{n+1} = g(x_n)
\end{equation}

can be implemented by storing the current approximation in a variable and
successively updating its value during a loop.

This model is closely related to the conventional state-oriented execution
model of modern processors and is naturally expressed in languages such as C.
Python also supports extensive use of imperative constructs, although it is a
multiparadigm language and provides additional abstractions.

Functional programming, in contrast, emphasizes function evaluation and
composition rather than explicit state modification. Common characteristics
include immutable values, first-class functions, higher-order functions,
function composition, and recursive definitions
\cite{hughes1989functional,hutton2016programming}.

In a functional formulation, the same numerical recurrence may be expressed
through repeated applications of a function:

\begin{equation}
    x_1 = g(x_0),
\end{equation}

\begin{equation}
    x_2 = g(g(x_0)),
\end{equation}

\begin{equation}
    x_3 = g(g(g(x_0))),
\end{equation}

and, more generally,

\begin{equation}
    x_n = g^{(n)}(x_0).
\end{equation}

This representation has a direct relationship with the mathematical
formulation of iterative methods, since computation may be expressed as
successive applications and compositions of functions rather than explicit
modifications of mutable state.

Functional programming also enables the use of higher-order functions, in
which functions may receive other functions as arguments or return them as
results. Hughes argues that higher-order functions and lazy evaluation can
serve as important mechanisms for improving modularity and decomposing programs
into reusable components
\cite{hughes1989functional}.

Haskell is one of the principal languages associated with purely functional
programming. The Haskell language is non-strict and purely functional, and
functions are first-class values. Its evaluation model allows expressions to
remain unevaluated until their values are required
\cite{marlow2010haskell,hutton2016programming}. The official Haskell 2010
Language Report defines the language and its non-strict semantics
\cite{marlow2010haskell}.

These properties enable algorithms to be expressed using abstractions that
differ substantially from conventional imperative implementations. At the
same time, lazy evaluation and immutable data structures may result in
execution and memory-access patterns that differ from those of strictly
evaluated imperative programs.